% Bagian: turing-machines
% Bab: machines-computations
% Subbagian: halting-states

\documentclass[../../../include/open-logic-section]{subfiles}

\begin{document}

\olfileid[id]{tur}{mac}{hal}
\olsection{Keadaan Penghentian}

\begin{explain}
Meskipun kita telah mendefinisikan bahwa mesin hanya berhenti ketika tidak ada
instruksi yang dapat dijalankan, representasi mesin Turing yang lazim memiliki
suatu \emph{keadaan penghentian} khusus~$h$, dengan $h \in Q$.

Gagasan di balik keadaan penghentian itu sederhana: ketika mesin telah selesai
beroperasi (siap menerima masukan atau telah selesai menuliskan keluaran),
mesin memasuki keadaan~$h$ dan berhenti di sana. Beberapa mesin memiliki dua
keadaan penghentian, satu yang menerima masukan dan satu yang menolak masukan.
\end{explain}

\begin{ex}\emph{Keadaan Penghentian}.
Untuk memperjelas konsep ini, mari kita mulai dengan memodifikasi mesin genap.
Alih-alih membiarkan mesin berhenti dalam keadaan~$q_0$ ketika masukannya
genap, kita dapat menambahkan instruksi yang mengirim mesin ke suatu keadaan
penghentian.
\[
\begin{tikzpicture}[->,>=stealth',shorten >=1pt,auto,node distance=2.8cm,
                    semithick]
  \tikzstyle{every state}=[fill=none,draw=black,text=black]

  \node[initial, state]         (A)                     {$q_0$};
  \node[state]         (B) [right of=A] {$q_1$};
  \node[state]         (C) [below of=A] {$h$};

  \path (A) edge [bend left] node {\TMtrans{\TMstroke}{\TMstroke}{R}} (B)
  	    edge node {\TMtrans{\TMblank}{\TMblank}{N}} (C)
        (B) edge [loop above] node {\TMtrans{\TMblank}{\TMblank}{R}} (B)
            edge [bend left] node {\TMtrans{\TMstroke}{\TMstroke}{R}} (A);
\end{tikzpicture}
\]

Mari kita perluas contoh ini lebih lanjut. Ketika mesin menentukan bahwa
masukannya ganjil, mesin tidak pernah berhenti. Kita dapat mengubah mesin agar
memiliki keadaan \emph{penolakan} dengan mengganti instruksi yang membentuk
kalang dengan instruksi untuk memasuki keadaan penolakan~$r$.
\[
\begin{tikzpicture}[->,>=stealth',shorten >=1pt,auto,node distance=2.8cm,
                    semithick]
  \tikzstyle{every state}=[fill=none,draw=black,text=black]

  \node[initial,state]         (A)                     {$q_0$};
  \node[state]         (B) [right of=A] {$q_1$};
  \node[state]         (C) [below of=A] {$h$};
  \node[state]         (D) [below of=B] {$r$};

  \path (A) edge [bend left] node {\TMtrans{\TMstroke}{\TMstroke}{R}} (B)
  	    edge node {\TMtrans{\TMblank}{\TMblank}{N}} (C)
        (B) edge node {\TMtrans{\TMblank}{\TMblank}{N}} (D)
            edge [bend left] node {\TMtrans{\TMstroke}{\TMstroke}{R}} (A);
\end{tikzpicture}
\]
\end{ex}

\begin{explain}
Penambahan keadaan penghentian khusus dapat bermanfaat dalam kasus seperti
ini karena memperlihatkan secara eksplisit kapan mesin menerima atau menolak
masukan tertentu. Namun, penting dicatat bahwa keadaan penghentian khusus
tidak menambah daya komputasi. Demikian pula, pengertian penghentian yang tidak
terlalu formal memiliki keunggulannya sendiri. Definisi penghentian yang
digunakan sejauh ini dalam bab ini membuat bukti \emph{Masalah Penghentian}
intuitif dan mudah diperagakan. Karena alasan itu, kita tetap menggunakan
definisi semula.
\end{explain}

\end{document}
